\section*{Capítulo \ref{ch:g8:solids} --- Pirâmides e cones}

\begin{solution}{exo:g8:solids:1}
Tetraedro: $4$ \omterm{def:g5:solids:def}{faces}, $6$ arestas, $4$ \omterm{def:g5:solids:def}{vértices} ($4 + 4 = 6 + 2$).
Base hexagonal: $7$ \omterm{def:g5:solids:def}{faces}, $12$ arestas, $7$ \omterm{def:g5:solids:def}{vértices}
($7 + 7 = 12 + 2$).
\end{solution}

\begin{solution}{exo:g8:solids:2}
$V = \frac13 \times 36 \times 10 = 120$ cm$^3$.

Base retangular: $B = 24$ cm$^2$, então
$V = \frac13 \times 24 \times 7.5 = 60$ cm$^3$.
\end{solution}

\begin{solution}{exo:g8:solids:3}
$V = \frac13 \pi \times 25 \times 9 = 75\pi \approx 236$ cm$^3$.
\end{solution}

\begin{solution}{exo:g8:solids:4}
\omterm{def:g7:areas:prism}{Cilindro}: $\pi \times 16 \times 12 = 192\pi \approx 603$ cm$^3$. \omterm{def:g8:solids:def}{Cone}:
$\frac{192\pi}{3} = 64\pi \approx 201$ cm$^3$. Razão: o \omterm{def:g8:solids:def}{cone} é um
terço do \omterm{def:g7:areas:prism}{cilindro}.
\end{solution}

\begin{solution}{exo:g8:solids:5}
$56 = \frac13 \times B \times 8$, então
$B = \frac{56 \times 3}{8} = 21$ cm$^2$.
\end{solution}

\begin{solution}{exo:g8:solids:6}
A superfície lateral se desenrola num setor de disco cujo \omterm{def:g3:shapes:shapes}{raio} é a
\emph{geratriz} $s$ (a distância do \omterm{def:g5:solids:def}{vértice} do topo até a borda ---
esse \omterm{def:g6:lines:objects}{segmento} fica sobre a superfície), e não a altura $h$.
\end{solution}

\begin{solution}{exo:g8:solids:7}
$V = \frac13 \times 35^2 \times 22 = \frac13 \times 1225 \times 22
= \frac{26\,950}{3} \approx 8\,983$ m$^3$ --- cerca de $9\,000$ metros
cúbicos.
\end{solution}

\begin{solution}{exo:g8:solids:8}
Altura: $h = \sqrt{13^2 - 5^2} = \sqrt{169 - 25} = \sqrt{144} = 12$
cm. \omterm{def:g6:measure:volume}{Volume}: $V = \frac13 \pi \times 25 \times 12 = 100\pi$ cm$^3$.
\end{solution}

\begin{solution}{exo:g8:solids:9}
O \omterm{def:g3:shapes:shapes}{raio} é o mesmo, então o \omterm{def:g6:measure:volume}{volume} do \omterm{def:g8:solids:def}{cone} é um terço do \omterm{def:g6:measure:volume}{volume} do
\omterm{def:g7:areas:prism}{cilindro} \emph{para a mesma altura}: o líquido alcança
$\frac{9}{3} = 3$ cm no \omterm{def:g7:areas:prism}{cilindro}.
\end{solution}

\begin{solution}{exo:g8:solids:10}
Areia $=$ um \omterm{def:g8:solids:def}{cone}: $V = \frac13 \pi \times 4 \times 4.5 = 6\pi
\approx 18.8$ cm$^3$. A ampulheta comporta dois \omterm{def:g8:solids:def}{cones} desses, então a
areia enche exatamente \omterm{def:g3:division:half}{metade} do \omterm{def:g6:measure:volume}{volume} interno total.
\end{solution}

\begin{solution}{exo:g8:solids:11}
\emph{1.} Sim: a fórmula $V = \frac13 Bh$ vale para qualquer posição
do \omterm{def:g5:solids:def}{vértice} do topo, desde que $h$ seja a distância perpendicular até o
plano da base. $V = \frac13 \times 36 \times 8 = 96$ cm$^3$.

\emph{2.} Diagonal da base: $\sqrt{6^2 + 6^2} = 6\sqrt2$ cm. A aresta
lateral mais longa é a hipotenusa de um \omterm{def:g6:shapes:triangles}{triângulo retângulo} de catetos
$6\sqrt2$ (no plano da base) e $8$ (vertical):
\[
\sqrt{(6\sqrt2)^2 + 8^2} = \sqrt{72 + 64} = \sqrt{136} \approx 11.7
\text{ cm}.
\]
\end{solution}

\begin{solution}{pb:g8:solids:1}
\textbf{1.} $G$ pertence à \omterm{def:g5:solids:def}{face} de cima $EFGH$ e às duas \omterm{def:g5:solids:def}{faces}
laterais $BCGF$ e $DCGH$. As três \omterm{def:g5:solids:def}{faces} que \emph{não} contêm $G$ são
o fundo $ABCD$, a frente $ABFE$ e a \omterm{def:g5:solids:def}{face} da esquerda $ADHE$ (todas
compartilham o \omterm{def:g5:solids:def}{vértice} $A$, o canto oposto a $G$). As três \omterm{def:g8:solids:def}{pirâmides}
são $ABCDG$, $ABFEG$ e $ADHEG$: cada uma tem uma \omterm{def:g5:solids:def}{face} quadrada do \omterm{def:g5:solids:def}{cubo}
como base e o \omterm{def:g5:solids:def}{vértice} distante $G$ no topo. Cada uma é uma \omterm{def:g8:solids:def}{pirâmide} de
base quadrada: $5$ \omterm{def:g5:solids:def}{faces}, $8$ arestas, $5$ \omterm{def:g5:solids:def}{vértices}
(\cref{ex:g8:solids:count}).

\textbf{2.} A diagonal longa $[AG]$ liga os dois únicos \omterm{def:g5:solids:def}{vértices} que
não pertencem a nenhuma das três bases. Uma rotação de um terço de
volta em torno da reta $(AG)$ leva o \omterm{def:g5:solids:def}{cubo} nele mesmo (as três arestas
que saem de $A$ --- em direção a $B$, $D$, $E$ --- são embaralhadas em
ciclo, e o mesmo acontece com as três arestas que chegam a $G$). Ela
leva a base $ABCD$ em $ADHE$, depois $ADHE$ em $ABFE$, e de volta:
\omterm{def:g5:solids:def}{vértice} do topo fixo em $G$, bases em ciclo. Cada \omterm{def:g8:solids:def}{pirâmide} é assim
levada na seguinte: as três são cópias idênticas.

\textbf{3.} Três peças idênticas que preenchem o \omterm{def:g5:solids:def}{cubo} repartem o
\omterm{def:g6:measure:volume}{volume} dele igualmente:
\[
V_{\text{pirâmide}} = \frac{a^3}{3} .
\]

\textbf{4.} A \omterm{def:g8:solids:def}{pirâmide} $ABCDG$ tem base $ABCD$, de \omterm{def:g6:measure:area}{área} $a^2$, e o
\omterm{def:g5:solids:def}{vértice} $G$ dela fica na vertical acima do canto $C$, à altura
$CG = a$ (uma aresta vertical do \omterm{def:g5:solids:def}{cubo}). Como no
\cref{exo:g8:solids:11}, a fórmula se aplica com a altura
perpendicular $h = a$:
\[
V = \frac13 \times a^2 \times a = \frac{a^3}{3} ,
\]
em acordo perfeito com a questão 3 --- a decomposição e a fórmula se
confirmam uma à outra.

\textbf{5.} Para $a = 6$: $V = \frac{6^3}{3} = \frac{216}{3} = 72$
cm$^3$ cada. A experiência de despejar água é essa decomposição em
forma líquida: o \omterm{def:g7:areas:prism}{prisma} sobre a base $ABCD$ com altura $a$ é o próprio
\omterm{def:g5:solids:def}{cubo}, e ele comporta exatamente três enchimentos da \omterm{def:g8:solids:def}{pirâmide} --- e
aqui os três enchimentos até se montam geometricamente no \omterm{def:g5:solids:def}{cubo}.

\textbf{6.} Ligar o centro $O$ aos quatro cantos de cada \omterm{def:g5:solids:def}{face} produz
seis \omterm{def:g8:solids:def}{pirâmides}, uma por \omterm{def:g5:solids:def}{face}: base quadrada, \omterm{def:g5:solids:def}{vértice} $O$ no topo. O
centro está igualmente posicionado em relação às seis \omterm{def:g5:solids:def}{faces} (qualquer
rotação ou simetria do \omterm{def:g5:solids:def}{cubo} fixa $O$ e permuta as \omterm{def:g5:solids:def}{faces}), então as seis
\omterm{def:g8:solids:def}{pirâmides} são idênticas. A altura de cada uma é a distância de $O$ até
uma \omterm{def:g5:solids:def}{face}: \omterm{def:g3:division:half}{metade} do lado, $\frac{a}{2}$.

\textbf{7.} Seis peças idênticas repartem o \omterm{def:g5:solids:def}{cubo}: $V = \frac{a^3}{6}$
cada.

\textbf{8.} Pela fórmula: \omterm{def:g6:measure:area}{área} da base $a^2$, altura $\frac a2$:
\[
V = \frac13 \times a^2 \times \frac{a}{2} = \frac{a^3}{6} .
\]
As duas respostas concordam.

\textbf{9.} As duas decomposições tratam de \omterm{def:g8:solids:def}{pirâmides} de proporções
bem especiais, talhadas num \omterm{def:g5:solids:def}{cubo}; uma \omterm{def:g8:solids:def}{pirâmide} esguia, uma \omterm{def:g8:solids:def}{pirâmide}
torta ou um \omterm{def:g8:solids:def}{cone} não podem ser montados assim (três \omterm{def:g8:solids:def}{cones} não enchem um
\omterm{def:g7:areas:prism}{cilindro}). É por isso que o \cref{thm:g8:solids:volume} é
\emph{admitido} neste nível: as decomposições e a experiência de
despejar água tornam o $\frac13$ inteiramente crível, e a demonstração
geral honesta --- a integração --- é dada no volume do ensino médio.

\textbf{10.} $V = \frac{6^3}{6} = 36$ cm$^3$; e, pela fórmula,
$V = \frac13 \times 36 \times 3 = 36$ cm$^3$. Concordam.

\textbf{11.} No triângulo $SAB$, os pontos médios dos lados $[SA]$ e
$[SB]$ são ligados por um \omterm{def:g6:lines:objects}{segmento} paralelo a $(AB)$ e de comprimento
$\frac{AB}{2} = \frac{c}{2}$ (\cref{thm:g8:midpoints:direct}). O mesmo
vale em cada \omterm{def:g5:solids:def}{face} lateral, então os quatro pontos médios formam um
quadrado de lado $\frac c2$, com os lados paralelos aos da base. Sobre
a reta vertical que passa pelo \omterm{def:g5:solids:def}{vértice} do topo, o corte passa pelo
ponto médio da altura (teorema dos pontos médios de novo, num
triângulo que passa por $S$, pelo centro da base e por um \omterm{def:g5:solids:def}{vértice} da
base): o plano de corte fica à altura $\frac h2$.

\textbf{12.} A peça de cima é uma \omterm{def:g8:solids:def}{pirâmide} regular de base quadrada de
lado $\frac c2$ e altura $\frac h2$ (a base dela é o corte, e o
\omterm{def:g5:solids:def}{vértice} do topo é $S$). O \omterm{def:g6:measure:volume}{volume} dela é
\[
\frac13 \times \left(\frac c2\right)^2 \times \frac h2
= \frac13 \times \frac{c^2}{4} \times \frac h2
= \frac{1}{8} \times \frac{c^2 h}{3}
= \frac{V}{8} :
\]
um oitavo do original --- e não a \omterm{def:g3:division:half}{metade}.

\textbf{13.} O tronco fica com o resto: $V - \frac V8 = \frac{7V}{8}$.
O corte na \omterm{def:g3:division:half}{metade} da altura reparte o \omterm{def:g6:measure:volume}{volume} na razão $1 : 7$ --- a
laje de baixo, de aparência modesta, guarda sete vezes o \omterm{def:g6:measure:volume}{volume} da
ponta de cima.

\textbf{14.} Abaixo da \omterm{def:g3:division:half}{metade} da altura fica $\frac78$ do \omterm{def:g6:measure:volume}{volume}. Com
$V = \frac13 \times 35^2 \times 22 = \frac{26\,950}{3} \approx
8\,983$ m$^3$ (\cref{exo:g8:solids:7}):
\[
\frac78 \times \frac{26\,950}{3} \approx 7\,860 \approx
7\,900 \text{ m}^3
\]
arredondando à centena --- a campanha da ``\omterm{def:g3:division:half}{metade} de baixo'' limpa
quase oito vezes o \omterm{def:g6:measure:volume}{volume} de vidro da outra.

\textbf{15.} Dobrar todas as dimensões multiplica a \omterm{def:g6:measure:area}{área} da base por
$2^2 = 4$ e a altura por $2$, logo o \omterm{def:g6:measure:volume}{volume} por
$4 \times 2 = 8 = 2^3$. Ampliar por $k$ multiplica as \omterm{def:g6:measure:area}{áreas} por $k^2$
e mais um comprimento por $k$: os \omterm{def:g6:measure:volume}{volumes} crescem pelo fator $k^3$. A
peça de cima da questão 12 é uma cópia reduzida da \omterm{def:g8:solids:def}{pirâmide} inteira com
$k = \frac12$, então o \omterm{def:g6:measure:volume}{volume} dela é
$\left(\frac12\right)^3 = \frac18$ do total --- a explicação em uma
linha.
\end{solution}
